% GENERATED FILE. Edit canonical lessons, then run npm run generate:books.
% canonical-source-hash: fnv1a64:b35827f48a21e898
% canonical-lessons: IT-C509-arbitro, IT-C509-fischiare, IT-C509-fallo, IT-C509-rigore, IT-C509-ammonire

\chapter{L'arbitro fischia}
\label{ch:it-l-arbitro-fischia}
\hlchaptermodality{509}

\begin{chapteropening}
\textbf{By the end of this chapter:} \emph{I can follow a referee's decisions in a match.}

Complete the last lesson of chapter 509: I can follow a referee's decisions in a match.
\end{chapteropening}

\section[l'arbitro]{l'arbitro — referee}

\label{lesson:IT-C509-arbitro}



\begin{quote}
Before the new one: say the Italian for abs, then the Italian for to sweat.
\end{quote}

\subsection*{You'll want to know: l'arbitro}
\textbf{l'arbitro} — \textquotedblleft{}referee\textquotedblright{}.

\begin{quote}
L'arbitro ha sbagliato tutto oggi.
\end{quote}

\begin{quote}
\emph{The referee got everything wrong today.}
\end{quote}

A woman is \emph{l'arbitra}.

\begin{cousinweb}
\emph{arbitro} is a learned borrowing from Latin \emph{arbiter} \textquotedblleft{}a witness, a judge\textquotedblright{} (Lewis \& Short).
\end{cousinweb}

\begin{grammarlens}[title={What you've built}]
One more word for telling, explaining and describing at B1.
\end{grammarlens}

\subsection*{Your turn}
\begin{itemize}
\raggedright
  \item \emph{Say it:} \emph{l'arbitro}
  \item \emph{Say it:} \emph{l'arbitro}, once more
  \item \emph{Say it:} a sentence of your own with \emph{l'arbitro}
  \item \emph{Recall:} say the Italian for abs, then the Italian for to sweat, then say \emph{l'arbitro} again
\end{itemize}

\begin{tcolorbox}[breakable,colback=teal!4,colframe=teal!35!black,title={Before you move on}]
What does l'arbitro mean? (\textquotedblleft{}Referee\textquotedblright{}.) Say it once more.
\end{tcolorbox}
\section[fischiare]{fischiare — to whistle; to boo}

\label{lesson:IT-C509-fischiare}



\begin{quote}
Before the new one: say the Italian for to sweat, then the Italian for referee.
\end{quote}

\subsection*{You'll want to know: fischiare}
\textbf{fischiare} — \textquotedblleft{}to whistle; to boo\textquotedblright{}.

\begin{quote}
Il pubblico ha fischiato la squadra.
\end{quote}

\begin{quote}
\emph{The crowd booed the team.}
\end{quote}

At a match, \emph{fischiare} is to boo: the opposite of \emph{applaudire}.

\begin{cousinweb}
\emph{fischiare} is inherited from Late Latin \emph{fistulāre} \textquotedblleft{}to play the pipe\textquotedblright{} (EtymDB), from \emph{fistula} \textquotedblleft{}a pipe\textquotedblright{} (Lewis \& Short).
\end{cousinweb}

\begin{grammarlens}[title={What you've built}]
One more word for telling, explaining and describing at B1.
\end{grammarlens}

\subsection*{Your turn}
\begin{itemize}
\raggedright
  \item \emph{Say it:} \emph{fischiare}
  \item \emph{Say it:} \emph{fischiare}, once more
  \item \emph{Say it:} a sentence of your own with \emph{fischiare}
  \item \emph{Recall:} say the Italian for to sweat, then the Italian for referee, then say \emph{fischiare} again
\end{itemize}

\begin{tcolorbox}[breakable,colback=teal!4,colframe=teal!35!black,title={Before you move on}]
What does fischiare mean? (\textquotedblleft{}To whistle; to boo\textquotedblright{}.) Say it once more.
\end{tcolorbox}
\section[il fallo]{il fallo — foul}

\label{lesson:IT-C509-fallo}



\begin{quote}
Before the new one: say the Italian for referee, then the Italian for to whistle.
\end{quote}

\subsection*{You'll want to know: il fallo}
\textbf{il fallo} — \textquotedblleft{}foul\textquotedblright{}.

\begin{quote}
Questo è fallo, arbitro!
\end{quote}

\begin{quote}
\emph{That is a foul, referee!}
\end{quote}

\emph{Fare fallo} — to commit a foul.

\begin{cousinweb}
\emph{fallo} \textquotedblleft{}fault, foul\textquotedblright{} comes from the old verb \emph{fallare} \textquotedblleft{}to err\textquotedblright{}, usually traced to Latin \emph{fallere} \textquotedblleft{}to deceive, to fail\textquotedblright{} (Lewis \& Short).
\end{cousinweb}

\begin{grammarlens}[title={What you've built}]
One more word for telling, explaining and describing at B1.
\end{grammarlens}

\subsection*{Your turn}
\begin{itemize}
\raggedright
  \item \emph{Say it:} \emph{il fallo}
  \item \emph{Say it:} \emph{il fallo}, once more
  \item \emph{Say it:} a sentence of your own with \emph{il fallo}
  \item \emph{Recall:} say the Italian for referee, then the Italian for to whistle, then say \emph{il fallo} again
\end{itemize}

\begin{tcolorbox}[breakable,colback=teal!4,colframe=teal!35!black,title={Before you move on}]
What does il fallo mean? (\textquotedblleft{}Foul\textquotedblright{}.) Say it once more.
\end{tcolorbox}
\section[il rigore]{il rigore — penalty (kick); strictness}

\label{lesson:IT-C509-rigore}



\begin{quote}
Before the new one: say the Italian for to whistle, then the Italian for foul.
\end{quote}

\subsection*{You'll want to know: il rigore}
\textbf{il rigore} — \textquotedblleft{}penalty (kick); strictness\textquotedblright{}.

\begin{quote}
L'arbitro ha dato il rigore all'ultimo minuto.
\end{quote}

\begin{quote}
\emph{The referee gave the penalty in the last minute.}
\end{quote}

Short for \emph{calcio di rigore}. Also strictness: \emph{con rigore}.

\begin{cousinweb}
\emph{rigore} is a learned borrowing from Latin \emph{rigor} \textquotedblleft{}stiffness, hardness\textquotedblright{} (Lewis \& Short); the penalty is the strict punishment of a foul.
\end{cousinweb}

\begin{grammarlens}[title={What you've built}]
One more word for telling, explaining and describing at B1.
\end{grammarlens}

\subsection*{Your turn}
\begin{itemize}
\raggedright
  \item \emph{Say it:} \emph{il rigore}
  \item \emph{Say it:} \emph{il rigore}, once more
  \item \emph{Say it:} a sentence of your own with \emph{il rigore}
  \item \emph{Recall:} say the Italian for to whistle, then the Italian for foul, then say \emph{il rigore} again
\end{itemize}

\begin{tcolorbox}[breakable,colback=teal!4,colframe=teal!35!black,title={Before you move on}]
What does il rigore mean? (\textquotedblleft{}Penalty (kick); strictness\textquotedblright{}.) Say it once more.
\end{tcolorbox}
\section[ammonire]{ammonire — to book, to caution; to warn}

\label{lesson:IT-C509-ammonire}



\begin{quote}
Before the new one: say the Italian for foul, then the Italian for penalty.
\end{quote}

\subsection*{You'll want to know: ammonire}
\textbf{ammonire} — \textquotedblleft{}to book, to caution; to warn\textquotedblright{}.

\begin{quote}
L'arbitro lo ha ammonito dopo il fallo.
\end{quote}

\begin{quote}
\emph{The referee booked him after the foul.}
\end{quote}

In football, a yellow card: \emph{è stato ammonito}.

\begin{cousinweb}
\emph{ammonire} is a learned borrowing from Latin \emph{admonēre} \textquotedblleft{}to remind, to warn\textquotedblright{} (EtymDB), from \emph{monēre} \textquotedblleft{}to warn\textquotedblright{}.
\end{cousinweb}

\begin{grammarlens}[title={What you've built}]
One more word for telling, explaining and describing at B1.
\end{grammarlens}

\subsection*{Your turn}
\begin{itemize}
\raggedright
  \item \emph{Say it:} \emph{ammonire}
  \item \emph{Say it:} \emph{ammonire}, once more
  \item \emph{Say it:} a sentence of your own with \emph{ammonire}
  \item \emph{Recall:} say the Italian for foul, then the Italian for penalty, then say \emph{ammonire} again
\end{itemize}

\begin{tcolorbox}[breakable,colback=teal!4,colframe=teal!35!black,title={Before you move on}]
What does ammonire mean? (\textquotedblleft{}To book, to caution; to warn\textquotedblright{}.) Say it once more.
\end{tcolorbox}
